% !TEX program = lualatex
% Transparencias de Fundamentos de las Matemáticas I
% Clase 09: Aplicaciones II

\documentclass[aspectratio=169,11pt]{beamer}

\usepackage{fontspec}
\usepackage[spanish,es-nodecimaldot]{babel}
\usepackage{amsmath,amssymb,mathtools}
\usepackage{booktabs}
\usepackage{csquotes}
\usepackage{xcolor}

\usetheme{Madrid}
\usecolortheme{default}
\setbeamertemplate{navigation symbols}{}
\setbeamertemplate{footline}[frame number]

\definecolor{FdMBlue}{RGB}{20,55,100}
\definecolor{FdMRed}{RGB}{130,30,30}
\definecolor{FdMGreen}{RGB}{25,100,60}

\setbeamercolor{structure}{fg=FdMBlue}
\setbeamercolor{title}{fg=white,bg=FdMBlue}
\setbeamercolor{frametitle}{fg=white,bg=FdMBlue}
\setbeamercolor{block title}{fg=white,bg=FdMBlue}
\setbeamercolor{block body}{bg=blue!3}
\setbeamercolor{alerted text}{fg=FdMRed}

\setmainfont{Latin Modern Roman}
\setsansfont{Latin Modern Sans}
\setmonofont{Latin Modern Mono}

\newcommand{\N}{\mathbb{N}}
\newcommand{\Z}{\mathbb{Z}}
\newcommand{\Q}{\mathbb{Q}}
\newcommand{\R}{\mathbb{R}}
\newcommand{\C}{\mathbb{C}}
\newcommand{\Pow}{\mathcal{P}}
\newcommand{\id}{\operatorname{id}}
\newcommand{\mcd}{\operatorname{mcd}}
\newcommand{\mcm}{\operatorname{mcm}}

\title[FdM I -- Clase 09]{Aplicaciones II}
\subtitle{Inyectividad, sobreyectividad, composición e inversa}
\author{Antonio Falcó Montesinos}
\institute{Universidad CEU Cardenal Herrera}
\date{Fundamentos de las Matemáticas I}

\begin{document}

\begin{frame}
  \titlepage
\end{frame}

\begin{frame}{Objetivos de la clase}
\begin{itemize}
  \item Caracterizar inyectividad, sobreyectividad y biyectividad.
  \item Trabajar con composición de aplicaciones.
  \item Entender cuándo existe una aplicación inversa.
\end{itemize}
\end{frame}

\begin{frame}{Mapa de la clase}
\tableofcontents
\end{frame}

\section{Tipos de aplicaciones}

\begin{frame}{Tipos de aplicaciones}
\begin{block}{Inyectividad}
\begin{itemize}
  \item \(f:A\to B\) es inyectiva si elementos distintos tienen imágenes distintas.
  \item Forma equivalente: si \(f(a_1)=f(a_2)\), entonces \(a_1=a_2\).
  \item Para probar inyectividad se suele partir de \(f(a_1)=f(a_2)\).
\end{itemize}
\end{block}
\end{frame}

\begin{frame}{Tipos de aplicaciones}
\begin{block}{Sobreyectividad}
\begin{itemize}
  \item \(f:A\to B\) es sobreyectiva si todo elemento de \(B\) es imagen de algún elemento de \(A\).
  \item Equivale a \(\operatorname{Im}(f)=B\).
  \item Para probarla, se toma \(b\in B\) y se busca \(a\in A\) tal que \(f(a)=b\).
\end{itemize}
\end{block}
\end{frame}

\begin{frame}{Tipos de aplicaciones}
\begin{block}{Biyectividad}
\begin{itemize}
  \item Una aplicación es biyectiva si es inyectiva y sobreyectiva.
  \item Cada elemento del codominio procede de un único elemento del dominio.
  \item Las biyecciones permiten construir inversas.
\end{itemize}
\end{block}
\end{frame}

\section{Operar con aplicaciones}

\begin{frame}{Operar con aplicaciones}
\begin{block}{Composición}
\begin{itemize}
  \item Si \(f:A\to B\) y \(g:B\to C\), entonces \(g\circ f:A\to C\).
  \item Se define por \((g\circ f)(a)=g(f(a))\).
  \item El orden importa: primero actúa \(f\), después \(g\).
\end{itemize}
\end{block}
\end{frame}

\begin{frame}{Operar con aplicaciones}
\begin{block}{Identidad e inversa}
\begin{itemize}
  \item \(\operatorname{id}_A:A\to A\) satisface \(\operatorname{id}_A(a)=a\).
  \item Si \(f:A\to B\) es biyectiva, existe \(f^{-1}:B\to A\).
  \item La inversa satisface \(f^{-1}\circ f=\operatorname{id}_A\) y \(f\circ f^{-1}=\operatorname{id}_B\).
\end{itemize}
\end{block}
\end{frame}

\section{Profundización}

\begin{frame}{Profundización}
\begin{block}{Pruebas tipo}
\begin{itemize}
  \item Inyectividad: partir de \(f(a_1)=f(a_2)\) y llegar a \(a_1=a_2\).
  \item Sobreyectividad: tomar \(b\in B\) y construir \(a\in A\) con \(f(a)=b\).
  \item Biyectividad: probar ambas propiedades o construir una inversa.
\end{itemize}
\end{block}
\end{frame}

\begin{frame}{Profundización}
\begin{block}{Composición: compatibilidad de dominios}
\begin{itemize}
  \item Para formar \(g\circ f\), la imagen de \(f\) debe estar dentro del dominio de \(g\).
  \item La notación \(g\circ f\) se lee de derecha a izquierda.
  \item No siempre existen \(g\circ f\) y \(f\circ g\) simultáneamente.
\end{itemize}
\end{block}
\end{frame}

\begin{frame}{Profundización}
\begin{block}{Inversa}
\begin{itemize}
  \item La inversa de una aplicación solo existe como aplicación si la original es biyectiva.
  \item La inversa deshace la acción de la aplicación.
  \item No debe confundirse \(f^{-1}(C)\) con una aplicación inversa.
\end{itemize}
\end{block}
\end{frame}

\begin{frame}{Cierre}
\begin{itemize}
  \item Inyectiva: no identifica elementos distintos.
  \item Sobreyectiva: alcanza todo el codominio.
  \item Biyectiva: permite invertir la asignación.
\end{itemize}
\end{frame}

\begin{frame}{Trabajo recomendado}
\begin{itemize}
  \item Releer las secciones correspondientes del manual.
  \item Identificar definiciones, símbolos nuevos y enunciados principales.
  \item Preparar dudas y ejemplos para el seminario asociado.
\end{itemize}
\end{frame}

\end{document}
